\documentclass[11pt]{article}

\usepackage[margin=1in]{geometry}
\usepackage{booktabs}
\usepackage{hyperref}
\usepackage{xcolor}
\usepackage{colortbl}
\usepackage{listings}
\usepackage{amsmath}
\usepackage{amssymb}
\usepackage{microtype}
\emergencystretch=3em

\hypersetup{
  colorlinks=true,
  linkcolor=blue,
  citecolor=blue,
  urlcolor=blue
}

\lstdefinestyle{code}{
  basicstyle=\ttfamily\small,
  breaklines=true,
  frame=single,
  columns=fullflexible,
  keepspaces=true,
  xleftmargin=0pt,
  xrightmargin=0pt
}
\lstset{style=code}

\newcommand{\impl}{\Rightarrow}
\newcommand{\TLA}{TLA\textsuperscript{+}}
\newcommand{\nxt}{\textsc{Next}}
\newcommand{\inv}{\textsc{Inv}}
\newcommand{\initp}{\textsc{Init}}
\newcommand{\prop}{\textsc{Property}}

\title{Proof-Driven Multi-Module Python via\\Agentic \TLA{} Co-Synthesis}
\author{Harry Wang \and Eric Shang Nan Chen}
\date{May 2026}

\begin{document}

\maketitle

\section{Introduction}

LLM-driven program synthesis has converged on a write-first, verify-second
loop: emit Python (or Dafny, or Rust), then run a test or a verifier and ask
the model to repair. This loop works on isolated functions, but the recent
DafnyComp benchmark~\cite{dafnycomp2025} shows that on compositional, multi-function
tasks, state-of-the-art LLMs hit only $3.69\%$ verified outputs --- a
$\approx 92$-point gap versus the single-function case. The bottleneck is not
syntax: DafnyComp reports $95.67\%$ syntactic correctness even on the
compositional split. The bottleneck is reasoning about how multiple modules,
each with its own invariants, fit together into a globally correct system.

\TLA{}~\cite{lamport2002specifying} is an attractive verification target for
that gap, for three reasons. First, \TLA{}'s refinement composition
pattern~\cite{hillelwayne-tla-adt} lets us discharge whole-system correctness
as a conjunction of per-child refinements, each verifiable by the same
model-checker (TLC~\cite{yu1999tlc}). Second, PlusCal~\cite{lamport2009pluscal}
gives LLMs an imperative algorithmic surface to write, which then translates
mechanically to \TLA{}'s state-transition style. Third, TLC's counterexamples
are structured state sequences that an LLM can read and revise --- the repair
signal is concrete rather than a binary pass/fail.

Our pipeline exploits all three properties in five stages: (i)~an LLM
\emph{decomposition} stage emitting a typed \texttt{DecompositionPlan} (parent +
up to four children); (ii)~\emph{bundled co-synthesis} of, per child, an abstract
spec module and an implementation module with an explicit abstraction map back to
parent variables; (iii)~a \emph{four-obligation verifier} (per-implementation
initiation, consecution, and property implication, plus a cross-module refinement
check) that pipes structured TLC counterexamples back to the synthesis agent;
(iv)~\emph{Python refinement} into one \texttt{icontract}-decorated~\cite{icontract}
class per module plus a composing parent app; and (v)~an advisory
\emph{trace-conformance gate} that replays one seeded Python execution against each
child's abstract spec via TLC~\cite{cirstea2024}. We evaluate on six compositional
benchmarks (producer--consumer, write-through cache, ledger audit, two-phase
commit, leader election, reader-writer lock) plus two legacy single-module
baselines, and add an instrumented efficiency study of the repair loop's token
cost and convergence.

\section{Background and Related Work}\label{sec:background}

\paragraph{\TLA{}, PlusCal, TLC.} \TLA{}~\cite{lamport2002specifying}
specifies systems as state machines using a temporal-logic formula
$\initp \wedge \Box [\nxt]_{vars}$. PlusCal~\cite{lamport2009pluscal} is an
algorithmic notation that compiles to \TLA{} (via \texttt{pcal.trans}). TLC~\cite{yu1999tlc}
is the explicit-state model checker for \TLA{} that we drive as a subprocess.

\paragraph{Refinement composition.} Our composition pattern follows Hillel
Wayne's \TLA{}-for-ADTs note~\cite{hillelwayne-tla-adt}: a parent module
declares variables and defines its own \initp{}/\nxt{}, and each child is
brought in as a named \texttt{INSTANCE} of its \emph{abstract} module via
\texttt{WITH} substitutions that bind the abstract variables to expressions
over the parent's variables (the abstraction map). The whole-system
refinement obligation is then $\bigwedge_i \texttt{Abs}_i!\texttt{Spec}$.
This pattern composes at the \TLA{} layer only --- PlusCal itself does not
compose --- so our pipeline glues children together after
\texttt{pcal.trans} has run on each implementation module.

\paragraph{DafnyComp.} DafnyComp~\cite{dafnycomp2025} is the September 2025
benchmark that motivates this work. It measures LLM-driven verified code
generation on both single-function and compositional multi-function tasks; the
collapse from $\approx 58\%$ to $3.69\%$ verification rate on the
compositional split is the headline gap our pipeline targets, from the \TLA{}
side rather than the Dafny side.

\paragraph{Trace conformance.} Cirstea et al.~\cite{cirstea2024} (\S 3.2)
encode a runtime trace as a literal sequence in \TLA{} and let TLC verify
that the trace is admitted by a spec. We adopt their encoding verbatim for
our Week-3 trace gate, with the modification that conformance is inferred from
the depth TLC reports against the trace length, since TLC 2.19 does not
implement a \texttt{POSTCONDITION} primitive.

\paragraph{Out of scope.} We do not use TLAPS for unbounded inductive proofs;
all our proofs are by finite model-checking with explicit small constant
domains. We do not use Nagini or Viper for Python verification; our Python
guarantee is the conjunction of (i)~static refinement and (ii)~runtime
\texttt{icontract} assertions plus the trace-conformance gate. We do not
attempt liveness; all four obligations are safety obligations.

\section{Approach}\label{sec:approach}

The pipeline has five stages (Figure~\ref{fig:pipeline}). All five are
wired through \texttt{run\_compositional\_pipeline} in
\texttt{src/main.py}; the legacy single-module path
(\texttt{run\_pipeline}) is preserved for the two non-compositional
baselines.

\begin{figure}[t]
\begin{lstlisting}
prompt --> PlannerAgent.decompose ----> DecompositionPlan
        -> SynthesisAgent.propose_bundle -> ModuleBundle
                                                |
            +-------------------------------------------+
            | Verifier.check_bundle:                    |
            |   per impl: pcal.trans + 3 TLC runs       |
            |   parent:   Refinement_<P>.tla + TLC      |
            +---------------+---------------------------+
                            |
              all_passed? no -> repair_bundle -> propose
                          yes
                            v
            RefineAgent.to_python_package -> PythonPackage
                            v
            TraceGate.check -> {child -> TraceResult}
                            v
            CompositionalPipelineResult(status, proof, traces)
\end{lstlisting}
\caption{The compositional pipeline. Arrows are typed data flow; the repair
loop is the only control-flow back-edge.}
\label{fig:pipeline}
\end{figure}

\subsection{Decomposition}
\texttt{PlannerAgent.decompose} (\texttt{src/agents/planner\_agent.py}) calls
Claude Opus 4.7 with a single function-calling tool,
\texttt{propose\_decomposition}, that emits a \texttt{DecompositionPlan}
consisting of a parent name plus 1--4 child modules. Each child carries an
\texttt{AbstractInterface} (state variables, action names, invariant sketch).
Decomposition is retried up to two times on Pydantic validation failure;
exhaustion raises \texttt{DecompositionError} and the pipeline returns
\texttt{status="planner\_failed"}.

\subsection{Bundled Co-Synthesis}
\texttt{SynthesisAgent.propose\_bundle} (\texttt{src/agents/synth\_agent.py})
takes the plan and asks the LLM, via \texttt{propose\_module\_bundle}, for a
\texttt{ModuleBundle}: a parent \TLA{} module plus, for each child, an
\emph{abstract} module and an \emph{implementation} module that carries
both PlusCal source and an \texttt{abstraction\_map} from abstract variables
to expressions over the parent's variables. The bundle schema (Pydantic
classes in \texttt{src/models/bundle.py}) enforces role consistency: only
implementations may carry PlusCal or an abstraction map; only one parent is
allowed.

\subsection{Four proof obligations}
Given a bundle the verifier (\texttt{src/agents/verifier.py}) discharges four
classes of obligation, summarised in Table~\ref{tbl:obligations}.

\begin{table}[t]
\centering
\small
\begin{tabular}{p{0.13\linewidth} p{0.21\linewidth} p{0.58\linewidth}}
\toprule
Obligation & Formula & TLC encoding (file) \\
\midrule
Initiation   & $\initp \impl \inv$ &
  \texttt{SPECIFICATION Spec, INVARIANT Inv} on
  \texttt{<Name>\_Impl.tla} (\texttt{tlc\_config.py:write\_init\_cfg}) \\
Consecution  & $\inv \wedge \nxt \impl \inv'$ &
  Aux \texttt{Consec\_<Name>\_Impl.tla} with
  \texttt{IInit == Inv}, \texttt{ISpec == IInit /\textbackslash{}
  [][Next]\_vars} (\texttt{tlc\_config.py:write\_consec\_module}) \\
Property impl. & $\inv \impl \prop$ &
  Aux \texttt{Prop\_<Name>\_Impl.tla} with
  \texttt{PInit == Inv}, unchanged transition
  (\texttt{tlc\_config.py:write\_property\_module}) \\
Refinement   & each \texttt{Impl} refines its \texttt{Abs} &
  Aux \texttt{Refinement\_<Parent>.tla}; see Listing~\ref{lst:refinement}
  (\texttt{formal/refinement.py:generate\_refinement\_module}) \\
\bottomrule
\end{tabular}
\caption{The four proof obligations and their TLC encodings.}
\label{tbl:obligations}
\end{table}

The first three obligations are per-implementation and reuse the
single-module encoding from the May 6 progress report. The fourth is the
compositional obligation: for the bundle as a whole, each child must refine
its abstract spec under the parent's substitution. We render this as the
Hillel-Wayne ADT pattern~\cite{hillelwayne-tla-adt}, shown in
Listing~\ref{lst:refinement}: the auxiliary module \emph{extends} the parent,
defines one named \texttt{INSTANCE} of the abstract module per child with
\texttt{WITH} substitutions for both the abstract variables and any abstract
constants, and asserts the conjunction of all \texttt{Abs}\textsubscript{i}!\texttt{Spec}
as a temporal property.

\begin{lstlisting}[caption={Generated \texttt{Refinement\_Parent.tla}
(simplified). The \texttt{const\_clauses} and \texttt{var\_clauses} come from
the abstract module's \texttt{CONSTANTS} declaration and the
implementation's \texttt{abstraction\_map}, respectively.},
label=lst:refinement,float=t]
---- MODULE Refinement_Parent ----
EXTENDS Parent
Abs_BackingStore == INSTANCE BackingStore_Abs WITH
    Keys <- Keys, Vals <- Vals,         (* CONSTANTS from abs *)
    storePresent <- store_p,            (* abstraction_map    *)
    storeVal     <- store_v
Abs_Cache == INSTANCE Cache_Abs WITH ...
RefinementSpec == /\ Abs_BackingStore!Spec
                  /\ Abs_Cache!Spec
====
\end{lstlisting}

TLC is then invoked on \texttt{Refinement\_<Parent>.tla} with
\texttt{SPECIFICATION Spec} (the parent's composed behaviour) and
\texttt{PROPERTY RefinementSpec}, which fails as soon as any child's
abstract behaviour is not admitted by its implementation under the
abstraction map.

\subsection{Counterexample-driven repair}
TLC's stdout is parsed by \texttt{src/formal/counterexample\_parser.py}
into an \texttt{ObligationResult} (one per failed check) carrying the
violated predicate and the offending state sequence. The synthesis agent's
\texttt{\_serialise\_bundle\_failure} flattens these per-implementation and
refinement failures into a single \texttt{tool\_result} block that the LLM
sees on its next turn. The system prompt is cached via Anthropic
ephemeral \texttt{cache\_control} so that repair iterations re-use the same
prompt without paying the input-token cost again.

\subsection{Refinement to \texttt{icontract} Python}
On a verified bundle, \texttt{RefineAgent.to\_python\_package}
(\texttt{src/agents/refine\_agent.py}) makes one LLM call per implementation
to emit a Python module --- one class per implementation, with
\texttt{@icontract.invariant} for each conjunct of the verified
\inv{}, and \texttt{@require}/\texttt{@ensure} for each method --- then one
more call to emit a parent \texttt{app.py} that composes the children. Each
mutating method calls a tiny \texttt{log\_action(name, state)} shim that
appends one JSONL line of (action name, abstract-variable snapshot) to a
trace file at runtime; this is the substrate for the trace gate.

\subsection{Trace conformance (advisory)}
\texttt{TraceGate.check} (\texttt{src/agents/trace\_gate.py}) subprocesses
the emitted package once with a seeded RNG and a step budget (default 50),
collects the \texttt{trace.jsonl} that the \texttt{log\_action} shim
produced, partitions entries by class prefix, and for each child renders a
\texttt{Trace\_<Name>.tla} module (Listing~\ref{lst:trace}) per
Cirstea et al.~\cite{cirstea2024}.

\begin{lstlisting}[caption={Generated \texttt{Trace\_<Name>.tla}, replaying
one execution against the child's abstract spec.},label=lst:trace,float=t]
---- MODULE Trace_Queue ----
EXTENDS Queue_Abs, Sequences, Naturals, TLC
Trace == << [event |-> "Enqueue", q |-> << >>],
            [event |-> "Enqueue", q |-> << 1 >>], ... >>
VARIABLE l
IsEvent(e) ==
    /\ l \in 1..(Len(Trace) - 1)
    /\ Trace[l + 1].event = e
    /\ q' = Trace[l + 1].q
    /\ l' = l + 1
TraceInit == l = 1 /\ q = Trace[1].q
TraceNext == \/ IsEvent("Enqueue") /\ Enqueue
             \/ IsEvent("Dequeue") /\ Dequeue
TraceSpec == TraceInit /\ [][TraceNext]_<<l, q>>
====
\end{lstlisting}

TLC is then run with \texttt{SPECIFICATION TraceSpec, INVARIANT Inv}. Since
TLC 2.19 has no \texttt{POSTCONDITION}, we classify the outcome from
TLC's reported depth: depth equal to the trace length with no error is
\texttt{conforms}; smaller depth with no error is \texttt{diverged at step
$K+1$}; an invariant violation is \texttt{invariant\_violated}; the rest is
\texttt{tlc\_timeout} or \texttt{tlc\_error}. Crucially, trace failures are
\emph{advisory only}: \texttt{CompositionalPipelineResult.status} is set by
the four static obligations, and trace results are reported alongside but do
not flip a verified bundle to unverified. The rationale is that the static
obligations already prove refinement over \emph{all} behaviours; the trace
gate is runtime evidence that the actual Python execution stays inside that
refinement.

\section{Implementation}\label{sec:impl}

The repository layout is shown in Figure~\ref{fig:layout}. The LLM backbone
is Claude Opus 4.7~\cite{anthropic-claude} with a three-tier overload
fallback: \texttt{claude-opus-4-7} $\xrightarrow{\texttt{529}}$
\texttt{claude-opus-4-6} $\xrightarrow{\texttt{529}}$ a tertiary OpenAI
adapter that translates between Anthropic-shape and OpenAI-shape
\texttt{tool\_use}. The verifier engine is TLC~2.19, driven as a subprocess
with \texttt{-config <obligation>.cfg -workers auto -deadlock}.

\begin{figure}[t]
\begin{lstlisting}
src/
  agents/   planner_agent, synth_agent, verifier,
            refine_agent, trace_gate
  formal/   pcal_translator, tla_runner, tlc_config,
            counterexample_parser, refinement, trace_renderer,
            trace_postcondition
  llm/      anthropic_client, openai_adapter, tools, prompts
  models/   task, decomposition, synthesis, bundle, proof
  main.py   run_pipeline + run_compositional_pipeline
  cli.py    Typer CLI (--compositional/--legacy)
examples/   8 benchmarks (2 legacy + 6 compositional)
scripts/    run_benchmarks.py (CSV + Markdown sweep)
tests/      27 test files, 221 offline cases + 1 live e2e
\end{lstlisting}
\caption{Repository layout.}
\label{fig:layout}
\end{figure}

The data plane is a single \texttt{CompositionalPipelineResult}
(\texttt{src/models/task.py}) carrying \texttt{status}, the
\texttt{DecompositionPlan}, the verified \texttt{ModuleBundle}, the
\texttt{CompositionalProofBundle} (one \texttt{ProofBundle} per child plus
the refinement \texttt{ObligationResult}), the emitted \texttt{PythonPackage},
the \texttt{traces} dict (one \texttt{TraceResult} per child), and a
\texttt{trace\_skipped\_reason} when the gate was not run (e.g.\ because the
bundle was unverified).

\section{Evaluation}\label{sec:eval}

We report three classes of evidence: an offline test suite that
exercises every component without an API key (\S\ref{ssec:tests});
the eight benchmark prompts that drive the live evaluation
(\S\ref{ssec:bench-catalog}), with per-benchmark detail in
\S\ref{ssec:bench-detail}; and the aggregate sweep results in
\S\ref{ssec:bench-agg}; an instrumented efficiency study of the repair
loop's cost and a controlled A/B of three convergence interventions
(\S\ref{ssec:reroll}); and a short qualitative postmortem of two structural
bugs closed before the sweep (\S\ref{ssec:qual}).

\subsection{Offline test suite}\label{ssec:tests}
The repository ships 27 \texttt{pytest} files exercising the pipeline at three
boundary levels: pure unit tests of the data models, counterexample parser,
trace renderer, and refinement-module generator; mocked-LLM tests of the planner,
synthesiser, and refiner that replay captured tool-call traces; and TLC-gated
integration tests that verify the trace gate against a real
\texttt{tla2tools.jar}. The offline suite runs at \textbf{221/221} passing without
an Anthropic API key, with external dependencies mocked at the boundary; a single
additional end-to-end test exercises the live API.

\subsection{Benchmark catalog}\label{ssec:bench-catalog}
The suite is two legacy single-module baselines plus six
compositional programs (Table~\ref{tbl:bench}). Each compositional
benchmark ships a \texttt{prompt.txt} (the natural-language requirement
sent to the planner) and a reference \texttt{expected\_modules.yaml}
(used only as evaluation diff --- the expected decomposition is
\emph{not} forced on the planner).

\begin{table}[t]
\centering
\small
\begin{tabular}{lll}
\toprule
Slug & Decomposition (parent + children) & Safety property \\
\midrule
\texttt{bounded\_counter}   & (single)                                & bounded in $0..10$ \\
\texttt{bank\_transfer}     & (single)                                & sum preserved \\
\midrule
\texttt{producer\_consumer} & Queue (abs+impl), Producer, Consumer & FIFO; queue $0..3$; producer bounded \\
\texttt{cache\_store}       & Store (abs+impl), Cache (abs+impl)   & read-after-write; cache${\subseteq}$store \\
\texttt{bank\_audit}        & Account, Ledger, Transfer            & conservation; ledger${\Leftrightarrow}$balance \\
\texttt{two\_phase\_commit} & ResourceManager, Coordinator         & agreement; no commit${\to}$abort \\
\texttt{leader\_log}        & Node, Election                       & ${\leq}1$ leader/term; logs monotone \\
\texttt{rw\_lock}           & RwLock, Client                       & mutex on write; reader count exact \\
\bottomrule
\end{tabular}
\caption{Benchmark suite. The six compositional benchmarks were
hand-authored to cover producer/consumer queues, write-through caches,
audit ledgers, distributed commit, leader election, and locking.}
\label{tbl:bench}
\end{table}

\subsection{Per-benchmark detail}\label{ssec:bench-detail}

We summarise each compositional benchmark by the decomposition the planner
converged on and its verification outcome; the prompts are summarised in
Table~\ref{tbl:bench}.

\paragraph{Verified on the first attempt (4 of 6).}
\texttt{cache\_store} (parent \texttt{CacheStoreSystem}; children \texttt{Store},
\texttt{Cache}) verified in 4 iterations. \texttt{two\_phase\_commit}
(\texttt{ResourceManager}, \texttt{Coordinator}) was fastest at 3 --- its small state
alphabet and explicit no-regression invariant make the inductive obligation easy.
\texttt{rw\_lock} (\texttt{RwLock}, \texttt{Client}) took 5, and \texttt{leader\_log}
(\texttt{Node}, \texttt{Election}) 6 --- the latter the highest TLC cost in the suite
(8.6\,s) from the 3-node $\times$ log-sequence product space. All per-child traces
conform.

\paragraph{Required a reroll (2 of 6).}
\texttt{producer\_consumer} (\texttt{Queue}, \texttt{Producer}, \texttt{Consumer})
exhausted its 6-iteration budget on the first attempt --- the \texttt{Queue} and
\texttt{Producer} impls kept failing per-module obligations while the LLM regenerated
structurally similar bundles --- then verified in 4 on a fresh reroll.
\texttt{bank\_audit} (\texttt{Account}, \texttt{Ledger}, \texttt{Transfer}) is the
hardest: its first two attempts passed every per-module obligation but failed the
cross-module \emph{refinement} check (the \texttt{INSTANCE WITH} bindings did not link),
and a third reroll verified in 5. Its conservation invariant couples the two accounts
through the ledger, so the abstraction map must project both balances and the ledger
onto each child.

\paragraph{Legacy baselines.}
\texttt{bounded\_counter} and \texttt{bank\_transfer} run through the single-module
\texttt{run\_pipeline} path (three obligations, no refinement) and both verify on the
first attempt; they are the May~6 progress-report baselines, preserved as regression
anchors.

\subsection{Aggregate results}\label{ssec:bench-agg}

Table~\ref{tbl:results} reports every CSV row produced by
\texttt{scripts/run\_benchmarks.py} across the initial sweep and the
two targeted reruns. Rows shaded in gray are the reroll attempts that
flipped a previously failing benchmark to verified.

\begin{table}[t]
\centering
\small
\setlength{\tabcolsep}{4pt}
\begin{tabular}{llcccccccc}
\toprule
Slug & Attempt & Status & It. & Ref. & Mod.\,P/F & Tr.\,c/d/iv & Pipe\,s & TLC\,s \\
\midrule
\texttt{cache\_store}       & initial  & verified   & 4 & passed & 2/0 & 2/0/0 & 201.3 & 6.6  \\
\texttt{leader\_log}        & initial  & verified   & 6 & passed & 2/0 & 2/0/0 & 387.2 & 8.6  \\
\texttt{rw\_lock}           & initial  & verified   & 5 & passed & 2/0 & 2/0/0 & 356.0 & 4.9  \\
\texttt{two\_phase\_commit} & initial  & verified   & 3 & passed & 2/0 & 2/0/0 & 182.7 & 5.2  \\
\texttt{producer\_consumer} & initial  & unverified & 6 & error  & 1/2 & 0/0/0 & 367.6 & 3.7  \\
\rowcolor{gray!12}
\texttt{producer\_consumer} & reroll-1 & \textbf{verified} & 4 & passed & 3/0 & 3/0/0 & 272.8 & 34.9 \\
\texttt{bank\_audit}        & initial  & unverified & 6 & error  & 2/0 & 0/0/0 & 430.0 & 5.1  \\
\texttt{bank\_audit}        & reroll-1 & unverified & 7 & error  & 2/0 & 0/0/0 & 358.2 & 5.5  \\
\rowcolor{gray!12}
\texttt{bank\_audit}        & reroll-2 & \textbf{verified} & 5 & passed & 2/0 & 2/0/0 & 544.2 & 5.5  \\
\bottomrule
\end{tabular}
\caption{Per-benchmark verification trajectories. Columns: repair
iteration count, refinement-obligation status, per-module passes/fails,
per-child trace outcomes (conform / diverged / invariant-violated),
total pipeline wall time, total TLC time. Shaded rows are the reroll
attempts that flipped a previously failing benchmark.
Source CSVs: \texttt{report/bench\_results.csv},
\texttt{report/bench\_rerun.csv}, \texttt{report/bench\_rerun\_2.csv}.}
\label{tbl:results}
\end{table}

\noindent The headline numbers are:

\begin{itemize}\itemsep0pt
  \item \textbf{Initial-attempt verification rate: 4 of 6 = 66.7\%.}
    Compared with DafnyComp's $3.69\%$ on
    compositional tasks~\cite{dafnycomp2025}, this is roughly a
    63-percentage-point improvement on the same problem class.
  \item \textbf{Cumulative verification rate after up to two
    rerolls: 6 of 6 = 100\%.}
  \item \textbf{Trace conformance is 100\% on every verified
    bundle.} All 13 per-child traces ($2+2+2+2+3+2$) conform; none
    diverge, none violate the invariant. The advisory trace gate
    fired on every verified run and stayed inside the proved
    refinement.
  \item \textbf{TLC time is a small fraction of total time}
    ($3.7$--$34.9$\,s of TLC vs.\ $182$--$544$\,s of wall clock).
    Pipeline cost is dominated by LLM round-trips, not model
    checking.
  \item \textbf{Iteration count does not predict success.}
    \texttt{two\_phase\_commit} converged in 3 iterations,
    \texttt{leader\_log} in 6; \texttt{producer\_consumer} failed at
    6 iterations on attempt one and then succeeded in 4 iterations
    on attempt two, with the same prompt and the same model.
\end{itemize}

\subsection{Efficiency: instrumentation and a controlled A/B}\label{ssec:reroll}

The sweep above suggested a tempting story: \texttt{producer\_consumer} failed at
its 6-iteration budget then verified in 4 on a fresh reroll, and
\texttt{bank\_audit} verified only on its third attempt --- as if a fresh start
``breaks an anchor'' in the repair conversation that more iterations cannot. We
built the machinery to test that story properly and to make the loop's cost
measurable. Three interventions were added, each behind a default-off flag so the
baseline is unchanged: (i)~\emph{history dedup} --- send the repair call only the
initial task plus the most recent failed attempt, instead of the full growing
transcript; (ii)~\emph{stuck-detection reroll} --- abandon a chain and re-propose
from scratch only when the same failure fingerprint repeats; and (iii)~a
deterministic \emph{pre-flight linter} that rejects reserved-PlusCal-label bundles
before spending a TLC run. We also instrumented every LLM call with its token usage,
tagged by pipeline stage.

The instrumentation gives the first quantified cost breakdown
(Table~\ref{tbl:cost}): \textbf{the \texttt{repair\_bundle} stage is \$2.07 of the
\$2.90 mean per-run spend} --- repair, not planning or refinement, is where the
money goes, confirming the premise that the repair loop is the dominant variable
cost.

\begin{table}[t]
\centering\small
\begin{tabular}{lcc}
\toprule
Stage & Mean \$/run & Mean input tokens \\
\midrule
planner       & 0.08 & 2473 \\
synth\_bundle & 0.35 & 1345 \\
\texttt{repair\_bundle} & \textbf{2.07} & \textbf{17531} \\
refine\_child & 0.30 & 3667 \\
refine\_parent& 0.24 & 4004 \\
\bottomrule
\end{tabular}
\caption{Per-stage cost (baseline condition, mean over 16 verified/unverified
runs). Estimated USD uses placeholder Opus prices; the \emph{relative} split is the
point. History dedup halves \texttt{repair\_bundle} input tokens (17.5k\,$\to$\,8.6k).}
\label{tbl:cost}
\end{table}

We then ran a controlled A/B: each lever in isolation, \texttt{-n 6}, three repeats
over all six benchmarks (Table~\ref{tbl:ab}). The result is deflating for the
anchoring story. \textbf{No condition is statistically distinguishable from baseline
on verification rate}: with 12--18 observations and per-benchmark rates swinging
0/3--3/3, binomial noise (SD\,$\approx$\,2 runs) swamps the 56--75\% spread. The
reroll ``win'' on \texttt{producer\_consumer} was a fresh coin-flip, not a reliable
escape --- automated stuck-reroll neither clearly helped nor hurt. History dedup
gave a \emph{modest, real} cost cut ($\sim$13\%), limited because Opus output tokens
dominate and dedup only trims input. Two caveats: the Anthropic credit balance was
exhausted mid-sweep, which invalidated one reroll run (so that arm is $N=2$); and
three runs crashed when the LLM emitted a malformed \texttt{ModuleBundle} that the
pipeline raised on rather than treating as a repairable failure --- a robustness gap
worth closing.

\begin{table}[t]
\centering\small
\begin{tabular}{lcccc}
\toprule
Condition & Verified & Rate & Mean iters & Mean \$/run \\
\midrule
OFF (baseline)                      & 12/18 & 67\% & 4.6 & 2.90 \\
DEDUP (latest-history)              & 10/18 & 56\% & 4.9 & 2.52 \\
REROLL-stuck ($N{=}2$)              & 9/12  & 75\% & 5.1 & 2.91 \\
\bottomrule
\end{tabular}
\caption{Controlled efficiency A/B (one lever each, default otherwise; the
credit-exhausted reroll run is excluded). The rate spread is within run-to-run
noise; dedup is the only consistent (modest) cost lever.}
\label{tbl:ab}
\end{table}

The honest takeaway is that the \emph{binding constraint is synthesis
nondeterminism, not repair-loop mechanics}: the same prompt and model verify
0/3--3/3, and reshaping the repair history or rerolling cannot move that. The
leverage lies upstream in synthesis quality (model, prompt, exemplars) or in simply
drawing more samples. The instrumentation --- localising cost to the repair stage
--- is the durable result; all three interventions ship default-off, and this A/B is
the evidence for leaving them so.

\subsection{Qualitative observations}\label{ssec:qual}
Two structural bugs found in the first end-to-end runs were closed
before the sweep that produced Table~\ref{tbl:results}. First, the
LLM repeatedly emitted PlusCal blocks whose labels collided with
PlusCal reserved names (\texttt{Done}, \texttt{Error},
\texttt{Lbl\_N}), causing \texttt{pcal.trans} to reject the module
before TLC was invoked --- now forbidden in
\texttt{SYNTH\_BUNDLE\_SYSTEM}. Second, the refinement aux module
generator initially omitted \texttt{CONSTANTS} substitutions in the
\texttt{INSTANCE} clauses, causing SANY to fail to link the parent's
abstract modules with cascading ``unknown operator'' errors --- now
fixed by parsing the abs module's \texttt{CONSTANTS} declaration and
emitting the matching \texttt{WITH} bindings
(\texttt{src/formal/refinement.py}). Both were observable through
the existing counterexample channel, but the LLM did not internalise
the mechanical rule from the error message alone. Lifting them into
hard prompt and code constraints unblocked the rest of the suite.

\section{Limitations and Future Work}\label{sec:limits}

\paragraph{Limitations.} The pipeline targets finite-domain inductive
checking with TLC; an unbounded inductive proof would need
TLAPS, which is out of scope. \inv{} must be expressible as an
initial-state predicate, so each variable must appear under explicit set
membership. The trace gate is per-child, against the child's abstract
spec; parent-level trace replay would require richer LLM coupling of
Python state to parent \TLA{} variables. Parametric abstract actions are
degraded: when an abstract operator takes arguments (e.g.\ \texttt{Enqueue(item)}),
the trace gate currently cannot reliably invent the parameter domain and
keeps only the per-state \inv{} check, dropping the action-relation
conjunct. Finally, the trace gate assumes the emitted Python is
deterministic: \texttt{random} is seeded at module import and the
subprocess runs under \texttt{PYTHONHASHSEED=0}; threading and unseeded
randomness are out of scope.

\paragraph{Future work.} Four follow-on items are queued.
(i)~\emph{Status split.} The single \texttt{status} field on
\texttt{CompositionalPipelineResult} splits into a triple of
\texttt{formal\_status}, \texttt{python\_status}, and
\texttt{overall\_status}, so the CLI's exit code reflects runtime
crashes in the emitted package, not just static verification. (ii)~\emph{Parametric-action fix.}
\texttt{log\_action} grows a \texttt{params} kwarg; the renderer reads it
and emits the action-relation conjunct, closing the parametric-action
degradation. (iii)~\emph{CrossHair gate.} \texttt{crosshair-tool} is added
as an optional gate behind a \texttt{--crosshair} flag, providing an
additional symbolic test of the emitted Python contracts.
(iv)~\emph{Enriched benchmark report.} \texttt{run\_benchmarks.py} is
extended with per-stage pass rates (TLC-runnable, invariant,
refinement, trace-conformance) plus action coverage and repair-iteration
count, so the headline DafnyComp-style comparison can be reported
directly.

\section{Conclusion}\label{sec:concl}

We have described an agentic \TLA{} pipeline that targets the
DafnyComp~\cite{dafnycomp2025} compositional-verification gap from the
\TLA{} side. The pipeline plans a parent-plus-children decomposition,
co-synthesises abstract and implementation modules for each child,
discharges three per-implementation TLC obligations plus a cross-module
refinement obligation in the Hillel-Wayne ADT style~\cite{hillelwayne-tla-adt},
emits an \texttt{icontract}-decorated Python package, and runs an advisory
trace-conformance gate~\cite{cirstea2024} against each child's abstract
spec. The four-obligation static proof is the load-bearing contribution;
the trace gate is a runtime witness that the emitted Python stays inside
the proved refinement. On the six-program suite the pipeline verifies
$\approx 67\%$ of bundles on the first attempt --- far above DafnyComp's
$3.69\%$ on the same problem class. An instrumented efficiency study localises
cost to the repair stage but finds that the binding constraint on convergence is
synthesis nondeterminism rather than repair-loop mechanics: reshaping the repair
history or rerolling does not move the verification rate beyond run-to-run noise.
The offline test suite is green at 221/221, and the convergence interventions ship
default-off, with the A/B as the evidence for leaving them so.

\bibliographystyle{plain}
\bibliography{refs}

\end{document}
